\scanpage{17}
\uline{Cas particuliers}

Supposons que $\forall i\in I$, \deleted{$E_i$ \unsure{soit|est} \unsure{rare|dense} dans $X$.} \inserted{on ait} (si $\Sigma_i^+\ \Sigma_i^-$ sont les deux comp. conn. de $X\setminus\Sigma_i$)

\deleted{Alors on \unsure{trouve} que}
\[
\begin{aligned}
\overline{\Sigma_i^+} &= \Sigma_i^+\cup E_i\\
\overline{\Sigma_i^-} &= \Sigma_i^-\cup E_i .
\end{aligned}
\]

Alors je dis qu'on a
\[
\overline{\Sigma_{12}}=\Sigma_{12}\cup E_1\cup E_2
\]

\noindent\rule{2.4cm}{0.3pt}\par

% slanted
\reversemarginpar
\margin{\unsure{introduire|introduisons} la partition $\Pi_j=\{\deleted{\unsure{\Sigma}}E_j,\Sigma_j^+,\Sigma_j^-\}$}
\normalmarginpar
\deleted{Supposons}

Soit $J\subset I$ tel que $I=J\amalg\{i\}$.

Considérons la partition \inserted{$\Pi_J$} de $X$ \deleted{\unsure{définie|définies|des parties}} définie par les $\deleted{\unsure{\Sigma}}E_j$ $(j\in J)$, \deleted{\unsure{et}} \unsure{intersection} des partitions \inserted{ternaires $\Pi_j$} de $X$ $(j\in J)$ définies par les $E_j$ \deleted{$\unsure{E|\in}$} $(j\in J)$
\[
X=E_j\amalg\coprod_{\alpha\in\pi_0(X\setminus E_j)}\Sigma_j^\alpha
\]

-- c'est aussi la partition définie par celle de $\displaystyle\bigcup_{j\in J}E_j$ en les $E_j$, et la partition \inserted{$\Pi_J^{\unsure{\circ|0}}$} de $\displaystyle X\setminus\mathop{\overline{\bigcup}}\limits_{j\in J}E_j$ intersection des partitions binaires $\Pi_{J,j}^{\circ}$ \inserted{$(j\in J)$} \unsure{induites|déduites} sur \inserted{$U_{\deleted{\unsure{J|j}}J}={}$} $\displaystyle X\setminus\bigcup_{j\in J}E_j$ par les $\Pi_j$ $(\unsure{g|j}\in J)$. P. \unsure{ex}
% archive sheet number: 16
